% ============================================================
% WORK_THREE: Análisis espectral FFT de animales, instrumentos y voces
% ============================================================

\subsection{WORK\_THREE: Análisis espectral FFT de sonidos de animales, instrumentos y voces}

Se realizó un análisis espectral mediante FFT (herramienta \texttt{signal\_lab}/\texttt{animal\_analyzer}) sobre ocho señales de audio: tres animales (gato, perro, gallo) y tres instrumentos (contrabajo, piano, flauta), todos audio real de dominio público/CC de Wikimedia Commons, y dos voces sintetizadas localmente con Windows SAPI (\texttt{David}/\texttt{Zira}) ante la falta de grabaciones humanas, documentadas como tales. Cada señal se cargó a su frecuencia de muestreo nativa (\texttt{librosa.load}, \texttt{sr=None}) y se acotó a una ventana fija de $3$~s. Para cada una se calculó su representación temporal y espectral (FFT de un lado, \texttt{numpy.fft.rfft}), su media, desviación estándar y frecuencia dominante $f_d$, y se introdujo un puntaje de separabilidad estadística entre pares, definido como la razón entre la diferencia de frecuencias dominantes y la suma de las desviaciones estándar de cada par ($S_{a,b} = \Delta f_d / (\sigma_a+\sigma_b)$, umbral $S_{a,b}\geq 3$ para separabilidad clara).

\begin{table}[!t]
\centering
\caption{Resultados del análisis de las señales (ventana de 3 s)}
\label{tab:wt-resultados}
\small
\begin{tabular}{@{}l l r r r@{}}
\toprule
Categoría & Señal & Media & Desv. Est. & $f_d$ (Hz) \\
\midrule
Animal & Gato & -0.0002 & 0.2508 & 1545.17 \\
Animal & Perro & -0.0000 & 0.1272 & 459.73 \\
Animal & Gallo & -0.0000 & 0.1708 & 1731.67 \\
\midrule
Instrumento & Contrabajo & 0.0000 & 0.0162 & 66.00 \\
Instrumento & Piano & 0.0000 & 0.1268 & 1179.33 \\
Instrumento & Flauta & 0.0000 & 0.1876 & 546.33 \\
\midrule
Voz$^{\dagger}$ & Persona 1 & 0.0196 & 0.1186 & 549.67 \\
Voz$^{\dagger}$ & Persona 2 & 0.0265 & 0.1213 & 180.00 \\
\bottomrule
\end{tabular}
\vspace{2pt}
{\footnotesize $^{\dagger}$Voz sintetizada (Windows SAPI), no grabación humana real.}
\end{table}

\begin{table}[!t]
\centering
\caption{Puntaje de separabilidad estadística por pares}
\label{tab:wt-separabilidad}
\small
\begin{tabular}{@{}l r r r l@{}}
\toprule
Par & $\Delta f_d$ (Hz) & $\sigma_a+\sigma_b$ & $S_{a,b}$ & Veredicto \\
\midrule
Gato--Perro & 1085.44 & 0.3781 & 2871.14 & Separables \\
Gato--Gallo & 186.50 & 0.4216 & 442.32 & Separables \\
Perro--Gallo & 1271.94 & 0.2981 & 4267.15 & Separables \\
\midrule
Contrabajo--Piano & 1113.33 & 0.1430 & 7787.68 & Separables \\
Contrabajo--Flauta & 480.33 & 0.2037 & 2357.60 & Separables \\
Piano--Flauta & 633.00 & 0.3143 & 2013.75 & Separables \\
\midrule
Persona 1--Persona 2 & 369.67 & 0.2399 & 1540.75 & Separables \\
\bottomrule
\end{tabular}
\end{table}

\begin{figure}[H]
\centering
\includegraphics[width=\linewidth]{work_three_animales_comparacion.png}
\caption{Comparación de los espectros frecuenciales de los sonidos animales.}
\label{fig:wt-animales}
\end{figure}

\begin{figure}[H]
\centering
\includegraphics[width=\linewidth]{work_three_instrumentos_comparacion.png}
\caption{Comparación de los espectros frecuenciales de los instrumentos musicales.}
\label{fig:wt-instrumentos}
\end{figure}

\begin{figure}[H]
\centering
\includegraphics[width=\linewidth]{work_three_personas_comparacion.png}
\caption{Comparación de los espectros frecuenciales de las dos voces (sintetizadas).}
\label{fig:wt-personas}
\end{figure}

Los siete pares evaluados resultaron, sin excepción, ``claramente separables'' según la métrica ($S_{a,b}$ entre $442$ y $7788$, muy por encima del umbral). Sin embargo, esta uniformidad expone una debilidad de diseño de la propia métrica: al dividir una diferencia de frecuencia dominante (en Hz, típicamente cientos o miles) por una suma de desviaciones estándar de amplitud normalizada por \texttt{librosa} (típicamente $\sigma \in [0.01, 0.3]$), el denominador es sistemáticamente pequeño para cualquier audio real, de modo que casi cualquier par con $f_d$ distintas supera con holgura el umbral. El puntaje pierde así poder discriminante entre pares ``muy separables'' y ``apenas separables''; se propone normalizar $\Delta f_d$ por el ancho de banda de Nyquist ($f_s/2$) en lugar de por $\sigma_a+\sigma_b$, o usar la dispersión de $f_d$ sobre subventanas, de modo que numerador y denominador compartan unidades de frecuencia.

Adicionalmente, dado que las dos señales de voz son sintetizadas (Windows SAPI) y no grabaciones humanas reales, sus resultados no deben interpretarse como evidencia sobre variabilidad de voces reales, sino solo como una comparación de pares reproducible con la misma herramienta. El reducido número de muestras (una grabación por señal) tampoco permite establecer conclusiones estadísticamente generalizables sobre las categorías estudiadas.
